\documentclass[../../hippo]{subfiles}
\begin{document}

\subsection{Derivation for Scaled Legendre (HiPPO-LegS)}
\label{sec:derivation-legs}

As discussed in \cref{sec:theory_legs},
the scaled Legendre is our only method that uses a measure with varying width.

\paragraph{Measure and Basis}
We instantiate the framework in the case
\begin{align}
  \omega(t, x) &= \frac{1}{t} \mathbb{I}_{[0, t]}
  \label{eq:scale-legendre-weight} \\
  g_n(t, x) &= p_n(t, x) = (2n+1)^{\frac{1}{2}} P_n\left( \frac{2x}{t} - 1 \right)
  \label{eq:scale-legendre-p}
\end{align}

Here, $P_n$ are the basic Legendre polynomials (\cref{sec:legendre-properties}).
We use no tilting, i.e.\ $\chi(t, x) = 1$, $\zeta(t) = 1$, and $\lambda_n=1$ so that
the functions $g_n(t, x)$ are an orthonormal basis.

\paragraph{Derivatives}

We first differentiate the measure and basis:
\begin{align*}
  \ppt \omega(t, \cdot)   &= -t^{-2}\mathbb{I}_{[0,t]} + t^{-1}\delta_t = t^{-1} (-\omega(t) + \delta_t)
  \nonumber \\
  \ppt g_n(t, x)
  &= -(2n+1)^{\frac{1}{2}} 2xt^{-2} P_n'\left( \frac{2x}{t} - 1 \right)
  \nonumber \\
  &= -(2n+1)^{\frac{1}{2}} t^{-1} \left(\frac{2x}{t} - 1 + 1\right) P_n'\left( \frac{2x}{t} - 1 \right)
  .
\end{align*}
Now define $z = \frac{2x}{t} - 1$ for shorthand and apply the properties of derivatives of Legendre polynomials (equation~\eqref{eq:legendre-xd}).
\begin{align*}
  \ppt g_n(t, x)
  &= -(2n+1)^{\frac{1}{2}} t^{-1} (z+1) P_n'\left( z \right)
  \\
  &= - (2n+1)^{\frac{1}{2}} t^{-1} \left[n P_n\left( z \right) + (2n-1)P_{n-1}(z) + (2n-3)P_{n-2}(z) + \dots \right]
  \nonumber \\
  &= -t^{-1} (2n+1)^{\frac{1}{2}} \left[ n (2n+1)^{-\frac{1}{2}} g_n(t, x) + (2n-1)^{\frac{1}{2}} g_{n-1}(t, x) + (2n-3)^{\frac{1}{2}}g_{n-2}(t, x) + \dots \right]
  \nonumber
\end{align*}

\paragraph{Coefficient Dynamics}

Plugging these into~\eqref{eq:coefficient-dynamics}, we obtain
\begin{align*}
  \frac{d}{dt} c_n(t)
  &= \int f(x) \left( \ppt g_n(t, x) \right) \omega(t, x) \dd x
  + \int f(x) g_n(t, x) \left( \ppt \omega(t, x) \right) \dd x
  \\
  &= -t^{-1} (2n+1)^{\frac{1}{2}}\left[ n(2n+1)^{-\frac{1}{2}}c_n(t) + (2n-1)^{\frac{1}{2}}c_{n-1}(t) + (2n-3)^{\frac{1}{2}}c_{n-2}(t) + \dots \right]
  \\
  &\qquad - t^{-1} c_n(t) + t^{-1} f(t)g_n(t,t)
  \\
  &= -t^{-1} (2n+1)^{\frac{1}{2}} \left[ (n+1)(2n+1)^{-\frac{1}{2}}c_n(t) + (2n-1)^{\frac{1}{2}} c_{n-1}(t) + (2n-3)^{\frac{1}{2}}c_{n-2}(t) + \dots \right]
  \\
  &\qquad + t^{-1} (2n+1)^{\frac{1}{2}} f(t) 
\end{align*}
where we have used $g_n(t, t) = (2n+1)^{\frac{1}{2}} P_n(1) = (2n+1)^{\frac{1}{2}}$.
Vectorizing this yields equation~\eqref{eq:scaled-legendre-dynamics}:
\begin{align}
  \ddt c(t) &= -\frac{1}{t} A c(t) + \frac{1}{t} B f(t)
  \label{eq:scaled-legendre-dynamics-details}
  \\
  A_{nk}
  &=
  \begin{cases}
    (2n+1)^{1/2}(2k+1)^{1/2} & \mbox{if } n > k \\
    n+1 & \mbox{if } n = k \\
    0 & \mbox{if } n < k
  \end{cases},
  \nonumber \\
  B_n &= (2n+1)^{\frac{1}{2}} \nonumber
\end{align}

Alternatively, we can write this as
\begin{equation}
  \ddt c(t) = -t^{-1} D
  \left[
    M D^{-1}
    c(t)
    +
    \mathbf{1}
    f(t)
  \right]
  ,
\end{equation}
where
$D := \diag \left[ (2n+1)^{\frac{1}{2}} \right]_{n=0}^{N-1}$,
$\mathbf{1}$ is the all ones vector, and the state matrix $M$ is
\begin{align*}%
  M
  =
  \begin{bmatrix}
    1 & 0 & 0 & 0 & \dots & 0 \\
    1 & 2 & 0 & 0 & \dots & 0 \\
    1 & 3 & 3 & 0 & \dots & 0 \\
    1 & 3 & 5 & 4 & \dots & 0 \\
    \vdots & \vdots & \vdots & \vdots & \ddots & \vdots \\
    1 & 3 & 5 & 7 & \dots & N \\
  \end{bmatrix},
  \quad \text{that is, } \quad
  M_{nk} =
  \begin{cases}
    2k+1 & \mbox{if } k < n \\
    k+1 & \mbox{if } k = n \\
    0 & \mbox{if } k > n
  \end{cases}
\end{align*}


Equation~\eqref{eq:scaled-legendre-dynamics-details} is a linear dynamical system, except dilated by a time-varying factor $t^{-1}$, which arises from the scaled measure.

\paragraph{Reconstruction}


By equation~\eqref{eq:reconstruction}, at every time $t$ we have
\begin{align*}
  f(x) \approx g^{(t)}(x)
  &= \sum_n c_n(t) g_n(t, x).
  \\
  &= \sum_n c_n(t) (2n+1)^{\frac{1}{2}} P_n\left( \frac{2x}{t} - 1 \right).
\end{align*}

\end{document}

